\chapter{Analysis from Ordering}
\label{chap:analysis-from-ordering}

% Closed question: How does a finite count of distinctions become an admissible
% notion of distance?

\begin{chapteressay}

The previous chapter introduced compatible records combining via
products: the algebra of records preserves source identity through
the operations. The chapter did not, however, introduce a notion of
distance between records. Two records of the same type might be
``close'' in some intuitive sense or ``far'' --- but the chapter's
algebra alone did not specify which.

This chapter takes up the question. It introduces distance as a
derived concept: not a primitive geometric notion, but a quantity
derived from the partial order on records. The chapter's claim is
that distance is forced by the order, not assumed separately. The
$p$-adic distance on integers is the canonical example: a metric
constructed from the divisibility relation alone, without any
reference to the integer line's geometry.

The chapter develops distance through five structures. The first is
distance from order: any partial order with a notion of refinement
gives rise to a metric. The $p$-adic distance is the canonical
example; Hamming distance and edit distance are others. Each metric
is a count of admissible distinctions separating two elements.

The second structure is the cubic spline. Given finite data, the
spline is the unique minimum-curvature interpolant. The Minimum
Curvature Property (Theorem 1) is the chapter's structural
result: among all admissible interpolants, the spline has the least
integrated squared curvature. This is Medium-section material: the
proof requires careful integration by parts and the natural
boundary conditions.

The third structure is the Galerkin projection. Given a continuous
target function and a finite-dimensional basis, the Galerkin
projection is the best approximation in the energy norm. The Best
Approximation Theorem is the chapter's second structural result;
the convergence theorem (fourth-order convergence for cubic
splines on $C^4$ targets) is the load-bearing technical content.
This is also Medium-section material.

The fourth structure is spline sufficiency. For $C^4$ targets,
cubic splines are sufficient for any practical resolution; below
that smoothness class, more specialized bases are needed. The
chapter's discipline is to choose the basis appropriate to the
target's smoothness.

The fifth structure is the Cauchy completion. Finite-dimensional
approximations have limits in the infinite-dimensional Sobolev
spaces. The completion guarantees the limit exists and is
admissible.

The chapter's ground example is the ratio of two counts, applied
to data points and their interpolation. Five data points (a small
counted record) are connected by a spline (a continuous
interpolant). The spline's shape is forced by the data and the
smoothness condition. The Galerkin projection generalizes to
target functions whose values are not all explicitly recorded; the
projection produces the best approximation in the basis space.

The chapter's second concrete example is lofting a boat hull. A
boatbuilder enlarges the designer's station drawings to full
size on the loft floor, then bends a long thin batten through the
station marks. The batten's natural shape is the chapter's spline
in physical form: a minimum-energy continuation of the data. The
loft lines are then transferred to patterns; the patterns are
used to cut the boat's wood. The discipline of the boatyard is
the discipline of the chapter: finite data forces an admissible
continuation through the spline; the continuation is
algorithmically realized through the Galerkin projection; the
physical result (the boat) is the chapter's structures performed
in wood.

The chapter's closed question is:

\begin{center}
\emph{How does a finite count of distinctions become an admissible
notion of distance?}
\end{center}

The chapter's answer is: distance is the count of admissible
distinctions, derived from the partial order; for continuous
structures, distance is the energy norm; the spline is the
minimum-curvature interpolant; the Galerkin projection is the
best approximation; the Cauchy completion provides the continuum.
The chain from finite data to continuous distance is forced
by the chapter's economy principle: only the minimum structure
required by the data is admissible.

This chapter and the next form the analytic core of the book. This
chapter establishes the spline and the Galerkin projection;
Chapter 6 builds on them to develop variational calculus and the
Euler-Lagrange equation. Together they give the mathematical gauge
its tools for selection and continuation.

\end{chapteressay}

\section{Distance from Order}
\label{se:distance-from-order}

The integers $12$ and $4$ differ by $8$. In the usual numerical sense,
the distance between them is $|12 - 4| = 8$. But the difference $8 =
2^3$ is divisible by $2$ three times. In the $2$-adic sense, the
integers $12$ and $4$ are close: they agree mod $2$, they agree mod
$4$, they agree mod $8$. They differ only at the next power of $2$.

The $p$-adic distance between two integers $a, b$ is defined as
$|a - b|_p = p^{-v}$, where $v$ is the largest power of $p$ that
divides $a - b$. For $p = 2$ and $a - b = 8 = 2^3$, the largest power
of $2$ dividing $8$ is $2^3$, so $v = 3$ and $|a - b|_2 = 2^{-3} =
1/8$. In the $2$-adic sense, $12$ and $4$ are close: distance $1/8$.

The $p$-adic distance is a genuine metric. It satisfies the triangle
inequality (in fact, the stronger ultrametric inequality $|x + y|_p \le
\max(|x|_p, |y|_p)$). It is symmetric ($|a - b|_p = |b - a|_p$). It
is positive definite ($|a - a|_p = 0$ and $|a - b|_p > 0$ when $a
\neq b$). Together these properties make $\mathbb{Z}$ a metric space
under the $p$-adic distance.

The $p$-adic distance is constructed entirely from the divisibility
order on the integers. There is no geometry, no number line, no notion
of ``how far apart on a ruler.'' The distance is a function of the
divisibility relation: two integers are close in the $p$-adic sense
when their difference is highly divisible by $p$. The order produces
the metric.

This is the chapter's first claim. Distance is not an independent
geometric concept; it can be derived from a partial order. The
divisibility lattice on $\mathbb{Z}$ has, for each prime $p$, an
associated metric. The metric is admissible because it satisfies
the metric axioms. The order produces the distance; the distance
does not exist separately from the order.

The chapter's metric examples extend beyond $p$-adics. The Hamming
distance between two strings of equal length is the number of
positions at which they differ. The string \texttt{ACGT} and the
string \texttt{ACCT} differ at one position; their Hamming distance
is $1$. The Hamming distance is derived from the order on positions:
two strings are close when most positions agree, far when many
differ.

The Levenshtein distance generalizes Hamming to strings of different
lengths: the minimum number of single-character insertions,
deletions, or substitutions that transform one string into the other.
Levenshtein is also order-based: the comparison is character-by-
character, with the cost computed from the alignment.

The chapter's structural claim: every metric on a discrete set is
fundamentally a counting metric: count the number of single-step
distinctions that separate two elements. The count is the distance.
The specific weights of the distinctions can vary (Hamming weights
each substitution equally; weighted Levenshtein weights insertions
and deletions differently), but the underlying structure is the
count of distinctions.

In Lean, the partial-order-based metric is the foundation of
\texttt{Spline.le} in \texttt{Episode6.lean}. The order on splines
produces a notion of distance: two splines are close when one is a
fine refinement of the other; far when their refinements diverge.
The distance is the chapter's count of admissible refinements
separating them.

\section{The Spline as Forced Continuation}
\label{se:the-spline-as-forced-continuation}

Five data points are recorded: $(x_0, y_0), (x_1, y_1), (x_2, y_2),
(x_3, y_3), (x_4, y_4)$ with $x_0 < x_1 < x_2 < x_3 < x_4$. The
question is: what function $f$ should connect these points? More
precisely: among all sufficiently smooth functions that pass through
every data point, which is the most economical --- the one introducing
the least unrecorded structure?

The answer is the natural cubic spline. The proof requires the
chapter's central theorem and its careful derivation.

The natural cubic spline through the five data points is a piecewise
cubic polynomial $S(x)$ defined on $[x_0, x_4]$ with the following
properties:
\begin{enumerate}
\item On each interval $[x_i, x_{i+1}]$, $S(x)$ is a cubic polynomial
$a_i + b_i (x - x_i) + c_i (x - x_i)^2 + d_i (x - x_i)^3$.
\item $S(x_i) = y_i$ for all $i$.
\item $S$ is twice continuously differentiable on $[x_0, x_4]$.
\item $S''(x_0) = S''(x_4) = 0$ (the natural boundary conditions).
\end{enumerate}

The first property says $S$ is piecewise cubic. The second says $S$
passes through every data point. The third says $S$, $S'$, and $S''$
are all continuous; this requires the pieces to match smoothly at the
interior knots. The fourth specifies the boundary behavior:
$S$ has zero curvature at the endpoints.

The theorem: the natural cubic spline through the five points exists
and is unique.

\emph{Proof of existence.} On each of the four intervals $[x_i,
x_{i+1}]$, $S$ has four coefficients $a_i, b_i, c_i, d_i$. The total
is $16$ unknowns. The constraints are:
\begin{itemize}
\item Five interpolation constraints: $S(x_i) = y_i$ for $i = 0, 1,
2, 3, 4$. But this is $5$ constraints; the first two intervals
share $S(x_1) = y_1$, etc. The count of independent interpolation
constraints is $5$ (each $y_i$ named once). Distributed across the
four intervals, this is $2$ constraints per interval ($S(x_i) =
y_i$ at the left endpoint, $S(x_{i+1}) = y_{i+1}$ at the right),
giving $8$ total constraints.
\item Continuity of $S'$ at the three interior knots $x_1, x_2,
x_3$: three constraints.
\item Continuity of $S''$ at the three interior knots: three
constraints.
\item Natural boundary conditions: $S''(x_0) = 0$ and $S''(x_4)
= 0$: two constraints.
\end{itemize}
Total: $8 + 3 + 3 + 2 = 16$ constraints. With $16$ unknowns and $16$
constraints, the system is square. Provided the constraints are
linearly independent (which can be shown for any distinct
$x_i$), the system has a unique solution. The solution exists.

\emph{Proof of uniqueness.} The system being linear with a unique
solution implies uniqueness; the spline is the unique cubic
polynomial-piecewise interpolant of the data satisfying the
smoothness conditions and the natural boundary.

Now the central economy theorem.

\begin{theorem}[Minimum Curvature Property]
Among all twice continuously differentiable functions $f$ on $[x_0,
x_4]$ that pass through the data points (i.e., $f(x_i) = y_i$ for
all $i$), the natural cubic spline $S$ minimizes the integrated
squared curvature:
\[
\int_{x_0}^{x_4} (f''(x))^2 \, dx \ge \int_{x_0}^{x_4} (S''(x))^2 \, dx
\]
with equality if and only if $f = S$.
\end{theorem}

\emph{Proof.} Let $g(x) = f(x) - S(x)$. Then $g(x_i) = 0$ for all $i$
(since $f$ and $S$ both pass through the data). Consider:
\[
\int_{x_0}^{x_4} (f'')^2 \, dx = \int_{x_0}^{x_4} (S'' + g'')^2 \, dx =
\int (S'')^2 \, dx + 2 \int S'' g'' \, dx + \int (g'')^2 \, dx.
\]

The middle term $\int S'' g'' \, dx$ can be integrated by parts twice:
\[
\int_{x_0}^{x_4} S'' g'' \, dx = [S'' g']_{x_0}^{x_4} - \int_{x_0}^{x_4}
S''' g' \, dx.
\]
The boundary term vanishes because $S''(x_0) = S''(x_4) = 0$ (natural
boundary). Integrating by parts again:
\[
- \int_{x_0}^{x_4} S''' g' \, dx = -[S''' g]_{x_0}^{x_4} + \int_{x_0}^{x_4}
S^{(4)} g \, dx.
\]
Since $S$ is piecewise cubic, $S^{(4)} = 0$ on each open interval. The
final term is $0$. The boundary term $[S''' g]$ also vanishes because
$g(x_0) = g(x_4) = 0$ (data interpolation forces this at the boundary).
Hence $\int S'' g'' \, dx = 0$.

Returning to the original expansion:
\[
\int (f'')^2 \, dx = \int (S'')^2 \, dx + \int (g'')^2 \, dx.
\]
Since $\int (g'')^2 \, dx \ge 0$ with equality iff $g'' = 0$
(which combined with $g(x_i) = 0$ forces $g \equiv 0$, hence $f = S$),
we conclude:
\[
\int (f'')^2 \, dx \ge \int (S'')^2 \, dx,
\]
with equality iff $f = S$. $\square$

The theorem establishes the spline's optimality property. Among all
admissible (twice-differentiable) interpolants of the data, the
spline has the minimum integrated squared curvature. The spline is
the smoothest interpolant in this precise sense.

The chapter's second claim is that the spline is forced. Given the
data and the smoothness condition (twice continuous
differentiability), the spline is not a choice; it is the unique
minimum-curvature interpolant. Any other interpolant would have more
curvature, which means more inflection between data points, which
means more unrecorded structure inserted into the gap.

The economy principle: a record should not assert structure beyond
what the data forces. The spline asserts only as much curvature as
is required to pass through the data while remaining
$C^2$-smooth. The minimum-curvature theorem is the precise
quantification of this economy.

In Lean, the \texttt{Spline} type in \texttt{Episode6.lean} carries
the spline's data: the knot positions, the coefficients on each
interval. The \texttt{Spline.le} relation orders splines by their
refinement: one spline refines another when it has additional knots
that the other lacks. The order is the chapter's distance from
order: two splines are close when their refinements differ by few
knot additions; far when they differ by many.

\begin{leanexcerpt}{Spline}
inductive Spline
  | observation : Prop $\to$ Spline
  | knot        : Prop $\to$ Polynomial $\to$ Spline $\to$ Spline
  | interpolant : Fact $\to$ Polynomial $\to$ Polynomial
                $\to$ Spline $\to$ Spline $\to$ Spline
\end{leanexcerpt}

The inductive carries the spline as a layered stack of refinements.
The \texttt{observation} constructor sits at the base: a single
propositional fact about the curve. The \texttt{knot} constructor
extends a prior spline by attaching a polynomial under a
propositional witness. The \texttt{interpolant} constructor pairs
two polynomials and two prior splines under a \texttt{Fact},
recording the interpolating piece that joins them. The chapter's
economy principle appears in the constructor count: only three
shapes are admissible, and a candidate spline must be expressible
as a finite tree of these constructors.

The natural boundary condition can be relaxed in other spline
variants. The clamped spline specifies $S'(x_0)$ and $S'(x_4)$
explicitly. The not-a-knot spline requires $S'''$ to be continuous
at $x_1$ and $x_{n-1}$ (essentially treating those knots as
non-knots). Each variant has its own boundary specification; each
has its own existence/uniqueness theorem. The natural spline is
the default; it minimizes curvature over the largest class of
admissible interpolants.

The chapter's economy theorem extends beyond cubic splines. For any
data and any chosen smoothness class, the minimum-norm interpolant
(minimum integrated squared $k$-th derivative for some $k$) is
admissible and unique. The cubic spline corresponds to $k = 2$
and natural boundary conditions; the quintic spline to $k = 3$;
and so on.

The discipline is: choose the lowest $k$ consistent with the data's
expected smoothness, and use the minimum-norm interpolant. Lower $k$
means less smoothness asserted; higher $k$ means more. The
appropriate choice depends on the source of the data: physical
measurements (typically smooth) suggest higher $k$; abstract
data without smoothness expectations suggest lower $k$.

\section{Galerkin Projection}
\label{se:galerkin-projection}

The previous section established the natural cubic spline as the
minimum-curvature interpolant of given data points. This section
takes the next step: given a continuous function $f$ on a bounded
interval, how does one find the spline that best approximates $f$ in
the energy norm? The answer is the Galerkin projection.

Let $V$ be the space of continuous functions on $[a, b]$ that vanish
at the endpoints (so $f(a) = f(b) = 0$). Equip $V$ with the inner
product
\[
\langle f, g \rangle = \int_a^b f'(x) g'(x) \, dx.
\]
This is an inner product (it is bilinear, symmetric, and positive
definite for non-constant functions). The induced norm $\|f\| =
\sqrt{\langle f, f \rangle} = \sqrt{\int (f')^2 \, dx}$ is the energy
norm or $H^1_0$-seminorm.

Now let $V_h \subset V$ be a finite-dimensional subspace: the
$n$-dimensional space spanned by $n$ basis functions $\phi_1, \ldots,
\phi_n$ chosen for the approximation. For cubic spline approximation,
the $\phi_i$ are the cubic B-splines centered at knots.

The Galerkin projection of $f$ onto $V_h$ is the element $f_h \in V_h$
satisfying:
\[
\langle f_h, v \rangle = \langle f, v \rangle \quad \text{for all } v
\in V_h.
\]
This is the variational characterization: $f_h$ is the element of
$V_h$ that has the same inner product with every test function $v
\in V_h$ as does $f$. Equivalently, the residual $f - f_h$ is
orthogonal to $V_h$ in the energy inner product.

The Galerkin projection can be computed by solving a linear system.
Write $f_h = \sum_j \alpha_j \phi_j$. The variational characterization
becomes:
\[
\sum_j \alpha_j \langle \phi_j, \phi_i \rangle = \langle f, \phi_i
\rangle \quad \text{for each } i = 1, \ldots, n.
\]
This is the system $K \alpha = F$ where $K_{ij} = \langle \phi_j,
\phi_i \rangle$ (the stiffness matrix), $F_i = \langle f, \phi_i
\rangle$ (the load vector), and $\alpha = (\alpha_1, \ldots,
\alpha_n)^T$ is the unknown coefficient vector. The matrix $K$ is
symmetric positive definite (because the inner product is); the
system has a unique solution $\alpha$; the Galerkin approximation
$f_h$ is constructed from $\alpha$.

\begin{theorem}[Galerkin Best Approximation]
Let $V_h \subset V$ be a finite-dimensional subspace. For any $f \in
V$, the Galerkin projection $f_h \in V_h$ is the best approximation
to $f$ in the energy norm:
\[
\|f - f_h\| = \min_{v \in V_h} \|f - v\|.
\]
\end{theorem}

\emph{Proof.} The variational characterization says $\langle f - f_h,
v \rangle = 0$ for all $v \in V_h$. For any $v \in V_h$, write
$v = f_h + (v - f_h)$ with $v - f_h \in V_h$. Then:
\[
\|f - v\|^2 = \|f - f_h - (v - f_h)\|^2.
\]
Expanding:
\[
= \|f - f_h\|^2 - 2 \langle f - f_h, v - f_h \rangle + \|v - f_h\|^2.
\]
The middle term vanishes by the variational characterization. So:
\[
\|f - v\|^2 = \|f - f_h\|^2 + \|v - f_h\|^2 \ge \|f - f_h\|^2,
\]
with equality iff $v = f_h$. Hence $f_h$ minimizes $\|f - v\|$ over
$V_h$. $\square$

This is the central convergence theorem. The Galerkin projection is
the closest element of $V_h$ to $f$ in the energy norm. As $V_h$
grows (more basis functions, finer knot spacing), the projection
$f_h$ improves: $\|f - f_h\|$ decreases monotonically.

The convergence rate depends on the smoothness of $f$ and the order
of the basis. For cubic B-splines with knot spacing $h$, the standard
result is:
\[
\|f - f_h\|_{L^2} \le C h^4 \|f^{(4)}\|_{L^2},
\]
provided $f \in H^4$ (four derivatives exist in $L^2$). The
exponent $4$ on $h$ is fourth-order convergence: halving the knot
spacing reduces the error by a factor of $16$. This is what makes
cubic splines such a powerful approximation tool.

The Galerkin projection is the chapter's algorithmic version of the
minimum-curvature theorem. The minimum-curvature theorem (previous
section) was a statement about the spline as the interpolant of
discrete data. The Galerkin projection is the same idea applied to
continuous functions: the spline space's best approximation to a
target function $f$ is the Galerkin projection. The two are
related: the minimum-curvature interpolant is the Galerkin
projection in a specific subspace.

In Lean, the Galerkin process is the typeclass
\texttt{GalerkinProcess} in \texttt{Episode6.lean}:

\begin{leanexcerpt}{GalerkinProcess}
structure GalerkinProcess
    (Value : Type)
    (Carrier : CarrierProcess Value)
    [DISTINGUISHABLE Value Carrier] [ADMISSIBLE Value Carrier]
    [COUNTABLE Value Carrier]       [ENCODED Value Carrier]
    [RESIDUE Value Carrier]         [BINARY Value Carrier]
    [REPEATABLE Value Carrier]      [NUMERIC Value Carrier]
    [REPRESENTABLE Value Carrier]   [PHYSICAL Value Carrier]
    [COMPARABLE Value Carrier]      [OBSERVED Value Carrier]
    [PRESENT Value Carrier]         [MEASURABLE Value Carrier]
    [GUNGAN Value Carrier]          [SOURCE Value Carrier]
    [EXECUTED Value Carrier]        [VALUE Value Carrier]
    [MAGNITUDE Value Carrier]       [SCALED Value Carrier]
    [LOAD Value Carrier]
  where
  ANSYS_process     : BASICProcess Value Carrier
  polynomial        : Polynomial
  scale_and_shift?  : Polynomial $\to$ Polynomial := ...
\end{leanexcerpt}

The structure carries the full cascade through \texttt{LOAD}, an
\texttt{ANSYS\_process} (a \texttt{BASICProcess} reference for the
basis), a current \texttt{polynomial}, and the function
\texttt{scale\_and\_shift?} that advances a polynomial by one
Galerkin step. The cascade is the proof obligation. An instance is
accepted only when every prior class in the stack has been
witnessed; the polynomial's progression then sits on certified
ground.

The \texttt{FINITE\_ELEPHANT} typeclass certifies that the
approximation is admissible at a stated resolution:

\begin{leanexcerpt}{FINITE_ELEPHANT}
class FINITE_ELEPHANT
    (Value : Type)
    (Carrier : CarrierProcess Value)
    [DISTINGUISHABLE Value Carrier] [ADMISSIBLE Value Carrier]
    [COUNTABLE Value Carrier]       [ENCODED Value Carrier]
    [RESIDUE Value Carrier]         [BINARY Value Carrier]
    [REPEATABLE Value Carrier]      [NUMERIC Value Carrier]
    [REPRESENTABLE Value Carrier]   [PHYSICAL Value Carrier]
    [COMPARABLE Value Carrier]      [OBSERVED Value Carrier]
    [PRESENT Value Carrier]         [MEASURABLE Value Carrier]
    [GUNGAN Value Carrier]          [SOURCE Value Carrier]
    [EXECUTED Value Carrier]        [VALUE Value Carrier]
    [MAGNITUDE Value Carrier]       [SCALED Value Carrier]
    [LOAD Value Carrier]
  where
  galerkin_process : GalerkinProcess Value Carrier
  finite?          : Polynomial $\to$ Polynomial $\to$ Prop := ...
\end{leanexcerpt}

The class binds the Galerkin process by a decision rule
\texttt{finite?} that takes two polynomials and returns the
proposition that one transforms to the other in finitely many
admissible steps. The cascade obligation runs the entire length of
the class stack; the decision rule covers the transition cases
(grounded constants, eigenvalue hits, Lanczos residual checks). The
compiler accepts the instance only when both the cascade and the
rule are provided.

The Galerkin projection extends beyond function approximation. It is
the foundation of the finite element method, the spectral method,
and many other numerical schemes. The general pattern: replace an
infinite-dimensional problem (find $f$ satisfying a differential
equation in $V$) with a finite-dimensional one (find $f_h$ in $V_h$
satisfying a discrete version of the equation). The Galerkin
projection guarantees that $f_h$ is the best $V_h$-approximation in
the energy norm.

For the differential equation $-u'' = g$ with boundary conditions
$u(a) = u(b) = 0$, the Galerkin formulation is: find $u_h \in V_h$
with $\langle u_h, v \rangle = \int g v \, dx$ for all $v \in V_h$.
The solution $u_h$ is the discrete approximation; as $V_h$ refines,
$u_h$ converges to the exact solution $u$.

The chapter's claim is that the Galerkin projection is the
algorithmic form of the chapter's economy principle. Among all
admissible discrete approximations to a continuous function, the
Galerkin projection minimizes the energy norm of the error. The
projection asserts only as much structure as the basis can
represent. The error is the residue; the residue's structure is
determined by the basis's limitations and the target's smoothness.

The Galerkin process is therefore the chapter's central
construction. It takes finite data (the basis, the target's inner
products with the basis) and produces a finite admissible
approximation. The approximation is the unique minimum-energy
element of the basis space matching the target's projections. The
construction is forced by the chapter's economy principle.

The minimum-curvature theorem of the previous section is a special
case: the data are the interpolation values at the knots; the
basis is the natural cubic spline space; the Galerkin projection
is the unique cubic spline interpolant. The general framework
includes data of many kinds (point evaluations, integrals,
boundary conditions) and bases of many forms (polynomials,
wavelets, eigenfunctions). The chapter's discipline operates
identically across all variations.

\section{Spline Sufficiency}
\label{se:spline-sufficiency}

The previous two sections established the natural cubic spline as the
minimum-curvature interpolant and the Galerkin projection as the best
approximation in the energy norm. This section asks: when is the
spline sufficient? When does the cubic spline approximation capture
enough of the target's structure that further refinement is not
needed?

The Spline Sufficiency property is the condition: for a target
function $f$ in the smoothness class $C^4$, the cubic spline
approximation with sufficiently fine knot spacing produces an error
below any stated tolerance. The fourth-order convergence theorem
quantifies this: $\|f - f_h\|_{L^2} \le C h^4 \|f^{(4)}\|_{L^2}$.

For most practical purposes, cubic splines are sufficient. The
fourth-order convergence is fast enough that modest knot spacings
($h$ of order $0.01$ to $0.001$) produce errors well below typical
measurement tolerances. The cubic spline is the workhorse of
practical interpolation and approximation.

The sufficiency condition fails for functions outside $C^4$. A
function with a discontinuity (a step function) cannot be
approximated by smooth splines without the Gibbs phenomenon
discussed in Chapter 3. A function with a corner (a kink in the
graph) has unbounded second derivative at the corner; cubic spline
approximation introduces a small oscillation near the corner.

For such non-smooth targets, higher-order or specialized
approximations are needed. Wavelet bases are well-suited to
functions with localized discontinuities. Piecewise constant or
piecewise linear approximations handle step functions. The choice
of basis is part of the calibration: the basis must be matched to
the target's smoothness class.

In Lean, the sufficiency condition is the proof obligation in
\texttt{FINITE\_ELEPHANT}. The instance must provide a bound on the
energy error; the bound is the resolution; the resolution is the
tolerance at which the approximation is admissible. If the target
is smoother than the basis can capture, the bound is tight; if the
target is rougher, the bound is looser.

The chapter's claim is that the spline is sufficient for the
chapter's purpose: providing an admissible discrete representation
of a continuous structure. For the chapter's questions about
distance, the spline's structure is enough to construct distances
on continuous functions. For more delicate questions (sharp
features, discontinuities, localized phenomena), additional
machinery is needed; the chapter does not develop it.

The chapter's discipline: the spline is sufficient when it is.
Outside its sufficiency class, the discipline is to acknowledge
the insufficiency and to seek a more appropriate basis. The
algorithmic form is the same (Galerkin projection); the basis
choice changes with the target's properties.

\section{The Cauchy Completion as Continuum}
\label{se:the-cauchy-completion-as-continuum}

Chapter 1 introduced the Cauchy completion of $\mathbb{Q}$ as the
construction of $\mathbb{R}$. This section completes the chapter by
extending the construction to the function spaces of interest.

The space of $L^2$ functions on $[a, b]$ is the Cauchy completion of
the space of continuous functions on $[a, b]$ under the $L^2$ norm:
\[
\|f\|_{L^2} = \sqrt{\int_a^b |f(x)|^2 \, dx}.
\]
A Cauchy sequence of continuous functions converges in $L^2$ to a
limit. The limit may not be continuous: it may have measure-zero
discontinuities. The $L^2$ space includes all such limits; it is
larger than the continuous functions.

The Sobolev space $H^1_0$ is the Cauchy completion under the energy
norm $\|f\| = \sqrt{\int (f')^2 \, dx}$. The functions in $H^1_0$ have
square-integrable first derivatives (in the weak sense) and vanish at
the boundary. The space is larger than the continuously
differentiable functions; it includes functions whose derivatives
exist only in a generalized sense.

The chapter's previous sections worked in the continuous functions or
in smooth subspaces; the Galerkin projection produces approximations
in finite-dimensional subspaces. The completion guarantees that the
approximation has a well-defined limit as the subspace refines. The
limit is in $L^2$ or in $H^1_0$, depending on which norm is used.

For cubic spline approximation of a $C^4$ target, the sequence
$f_{h_n}$ with $h_n \to 0$ is Cauchy in both $L^2$ and $H^1_0$ (in
fact in any Sobolev norm up to $H^2$, by the standard convergence
theorem). The limit is the target $f$ itself. The Cauchy completion
is needed because the approximations are in finite-dimensional
subspaces, but the target lives in the infinite-dimensional space
of continuous functions.

The chapter's fifth claim: the Cauchy completion of the discrete
approximations is the continuum. The continuum is not assumed; it
is constructed as the completion of finite-dimensional
approximations. The construction is constructive: each step
produces a finite-dimensional approximation; the limit is the
continuum's element.

In Lean, the completion is part of \texttt{Mathlib}'s functional-
analysis library. The type \texttt{Lp} represents $L^p$ functions
on a measure space. The Sobolev spaces $H^k$ are defined as the
completion of smooth functions under the corresponding Sobolev
norms. The chapter's algebraic structures (spline,
\texttt{GalerkinProcess}, \texttt{FINITE\_ELEPHANT}) operate in
these completed spaces.

The chapter closes here. Distance from order is a metric derived
from a partial order. The cubic spline is the minimum-curvature
interpolant of finite data. The Galerkin projection is the best
approximation in a finite-dimensional subspace. The spline is
sufficient for $C^4$ targets. The Cauchy completion provides the
continuum's underlying space.

These five structures together produce the chapter's answer to its
closed question. A finite count of distinctions becomes an
admissible notion of distance through the partial-order metric;
the metric extends to function spaces via the Galerkin projection;
the projection produces an admissible approximation; the
approximation's error is bounded by the resolution; the
approximation's limit, as the resolution refines, is the
continuum's element.

\begin{bridge}

The chapter's question was how a finite count of distinctions
becomes an admissible notion of distance.

The answer is: through the partial-order metric and the Galerkin
projection. The $p$-adic distance is the canonical example of a
metric derived from a partial order on integers. The Hamming
distance on strings and the divisibility distance on integers are
other examples; each is a count of admissible distinctions. The
cubic spline is the minimum-curvature interpolant of finite data;
its uniqueness theorem (Minimum Curvature Property) is the
chapter's structural result. The Galerkin projection is the best
approximation in the energy norm; the Best Approximation theorem
guarantees that the projection is the closest element of the
finite-dimensional subspace to the target. Spline sufficiency
identifies the smoothness class in which cubic splines are
adequate. The Cauchy completion provides the underlying continuum
in which the approximations live.

What remains is the selection of one continuation from many. The
spline is the minimum-curvature continuation; the Galerkin
projection is the minimum-error approximation; but in some
situations, the data underdetermines the answer, and many
admissible continuations exist. Chapter 6 takes up the question:
what selects one admissible continuation as the history actually
carried forward? The chapter introduces variational calculus and
its Euler-Lagrange equation as the structural mechanism for
selection.

\end{bridge}

\begin{coda}{Lofting a Boat Hull}

The lofting floor is in the second-story workshop above the boatyard.
Its surface is twenty-four meters long and six meters wide, planked
in straight-grained Douglas fir, sanded smooth, painted matte black.
This is the loft. On its surface, full-size drawings of the boat's
shapes will be made: each frame, each plank, each beam will be drawn
at full scale before any wood is cut.

The boatbuilder has the designer's plans: a small set of cross-section
drawings showing the boat's hull at twelve stations along the keel. The
stations are spaced at uniform intervals. Each station drawing shows
the hull's outline at that point: the half-breadth (width from the
centerline), the height of the keel, the height of the deck, the
curve of the hull. The drawings are the data; the boatbuilder's task
is to construct the full hull from them.

The first task is to enlarge the station drawings to full size. The
designer's drawings are at a scale of $1:10$; the loft drawing is at
$1:1$. The boatbuilder uses a proportional divider to step off the
designer's measurements at ten times the scale, marking the loft
floor with chalk lines and small nails. After several hours, twelve
station outlines are drawn on the loft, each at the location
corresponding to its station number.

The station drawings are the data points of the chapter's spline. Each
station provides values for the hull's shape at one specific location
along the keel. Between stations, the hull's shape is not directly
specified; the boatbuilder must fill in the gap.

Now comes the lofting proper. The boatbuilder takes a long thin
batten --- a flexible strip of wood, perhaps four meters long and a
centimeter wide --- and bends it gently along the marks of the station
outlines. The batten passes through the marked points; between them,
it takes its natural curved shape. The natural curve is the
minimum-energy continuation of the data points.

This is the chapter's spline in physical form. The batten is the
spline; the marks at the stations are the data points; the natural
curve of the batten between marks is the cubic-like continuation;
the batten's natural shape minimizes the integrated curvature in the
mathematically precise sense (within the elastic limits of the
wood).

The boatbuilder verifies the curve. She sights along the batten from
one end. Any bump, any flat spot, any unnatural inflection would
indicate that the batten is not in its minimum-energy state: that
the data points are inconsistent with the smooth hull they should
produce, or that the batten is being constrained somewhere it
should not be. A bump suggests a measurement error at a station;
a flat spot suggests two stations should have intermediate
information; an unnatural inflection suggests a structural problem
with the hull's design.

For this hull, the verification succeeds. The batten bends smoothly
through all twelve station marks without bumps or flat spots. The
hull's full longitudinal curve is now visible on the loft floor.

But the batten is one curve. The hull has many curves: the sheer
line (the upper edge of the side), the chine (where the side meets
the bottom), the keel (the bottom-most curve), and the diagonals.
Each is its own spline; each is drawn separately on the loft. The
spline structure of the chapter applies independently to each curve:
each is a minimum-energy continuation of the station data; each
admits a Galerkin projection; each is sufficient for $C^4$ targets.

The boat's hull is therefore the product of multiple splines. The
product structure preserves each spline's identity: the sheer line
is a sheer line, not just an unidentified curve; the chine is a
chine; the keel is a keel. Each is admissible separately; the hull
is the admissible composition.

After all the lines are drawn on the loft, the boatbuilder uses
them to make patterns: full-size templates for each plank, each
frame, each beam. The patterns are made by transferring the loft
lines to thin plywood. The plywood patterns are then used as
guides for cutting the actual hull wood.

The transfer is the algorithmic version of the Galerkin projection.
The continuous loft lines (an infinite-dimensional curve in
principle) are projected onto the discrete pattern surfaces (a
finite-dimensional approximation). The projection is the
boatbuilder's tracing: a careful sequence of measurements and
copies that produces the pattern's shape. The Galerkin best-
approximation theorem guarantees that the pattern is the closest
admissible approximation to the loft line, given the constraints
of the pattern material.

The patterns are then taken to the boatyard floor. The wood is laid
out flat; the patterns are placed on top; the shapes are marked.
The marks are cut. The cut wood is bent into the boat's frames and
attached to the keel. The boat takes shape physically.

Throughout this process, the chapter's discipline operates. The
spline is the unique minimum-curvature continuation of the station
data. The Galerkin projection is the best discrete approximation
to the continuous spline. The patterns are the finite-dimensional
realization of the projection. The cut wood is the physical
embodiment of the patterns. None of the operations introduce
unrecorded structure: the boat's hull is forced by the station
data through the chapter's algebraic operations.

When the boat is launched, the hull's shape is determined by the
station data and the spline's economy principle. A different
designer might have used different station data; the resulting
hull would differ. A different lofting technique would produce
different finite-dimensional approximations; the resulting hulls
would converge to the same continuous curve as the resolution
refines.

The boatyard's discipline is the chapter's discipline. The
station data are the recorded marks. The spline is the admissible
continuation. The Galerkin projection is the admissible
approximation. The patterns are the discrete realization. The
hull is the physical product. The boat floats because the algebra
is correct.

This is the chapter's claim in maritime form. Finite data, admissible
continuation, minimum-curvature interpolation, Galerkin projection,
finite-dimensional approximation, physical realization --- all of
these are the chapter's structures, performed in the boatyard. The
mathematics is not separate from the craft; the mathematics is the
craft's algebraic content.

The hull launches in three months. The boatbuilder will return to
the loft tomorrow to draw the next set of patterns. The discipline
continues; the algebra continues; the boat takes shape one plank at
a time, each plank forced by the loft lines, each loft line forced
by the station data, each station data point recorded with the
care that the chapter's discipline requires.

\end{coda}
